\documentclass[11pt]{article}
\usepackage{amsmath,amssymb,amsthm,mathtools}
\usepackage[margin=1in]{geometry}
\theoremstyle{plain}
\newtheorem{lemma}{Lemma}
\theoremstyle{remark}
\newtheorem{remark}{Remark}
\newcommand{\Z}{\mathbb{Z}}\newcommand{\E}{\mathbb{E}}
\newcommand{\bone}{\mathbf 1}

\title{Formalization brief: \texttt{lem:crossbridge} on a finite lattice}
\date{4 July 2026}

\begin{document}
\maketitle

\noindent\textbf{Goal.} Machine-check the algebraic core of \texttt{lem:crossbridge}
(the dual pair computes the species cross-correlation) on a \underline{finite} lattice,
in a new self-contained file (suggest \texttt{TypeDDecouplingCrossbridge.lean}). The
continuum statement in \texttt{TypeDDecouplingCrossover.lean} stays as is (its
\texttt{sorry} covers the scaling embedding); add only a docstring cross-reference
recording that its finite-lattice core is now proved.

\section*{Setting}

Finite lattice $\Lambda=\{-L,\dots,L\}$, two-species type D ASEP at $n=\infty$
(the rate table of \texttt{eq:rates}, already encoded in the project), generator
$\mathcal L=\sum_b\mathcal L_b$ over bonds, all as finite matrices on the
configuration space. The \underline{triangular} duality function, for a dual configuration
$\xi$ with finitely many particles and a process configuration $\eta$:
\[
  D^{\mathrm{tri}}(\xi,\eta)=\bone\{\xi\subseteq\eta\}
  \prod_{i=1,2}\prod_{x:\,\xi_{i,x}=1}
  q^{2\left(x-N^-_{x-1}(\xi_i)+N^+_{x+1}(\eta_i)\right)} .
\]
For the crossbridge only the \underline{two-particle sector} is needed: $\xi$ has exactly
one particle of each species (positions $(x_1,x_2)$), and the dual dynamics on that
sector is the two-particle process already encoded (the coupled pair with the
$n=\infty$ rates; cf.\ \texttt{TypeDDecouplingDualPairWitness.lean}).

\section*{Step 1: interlacing $\Rightarrow$ semigroup duality (finite linear algebra)}

If $\mathcal L\,D=D\,\mathcal L_{\mathrm{dual}}^{\mathsf T}$ as finite matrices
($D$ the two-particle-sector duality matrix $D[\xi,\eta]:=D^{\mathrm{tri}}(\xi,\eta)$),
then for all $s\ge0$, $e^{s\mathcal L}D=D\,e^{s\mathcal L_{\mathrm{dual}}^{\mathsf T}}$:
from $AD=DB$ one gets $A^kD=DB^k$ by induction and the exponential series converges
entrywise (finite dimension); Mathlib's \texttt{Matrix.exp} suffices. This gives the
duality relation
\begin{equation}\label{eq:dualrel}
  \E_\eta\big[D^{\mathrm{tri}}(\xi,\eta_s)\big]
  =\E_\xi\big[D^{\mathrm{tri}}(\xi_s,\eta)\big].
\end{equation}

\underline{Primary goal} for the input: prove the two-particle-sector interlacing
$\mathcal L D=D\mathcal L_{\mathrm{dual}}^{\mathsf T}$ itself, bond by bond: for each
bond $b$, the identity $\mathcal L_b D=D(\mathcal L^{\mathrm{dual}}_b)^{\mathsf T}$
is finite algebra --- the exponent bookkeeping under a bond-local move is explicit
($N^+_{x+1}(\eta_i)$ changes only when the move crosses $x$, and the rate table's
$1/q^2/\eps$ structure matches the $q^{2}$-shifts; this is the type D instance of
the standard telescoping check, and in the two-particle sector the matrices are
small). \underline{Sanctioned fallback}: if the general-$L$ bondwise identity resists,
take the two-particle-sector interlacing as a \emph{single named hypothesis} (this
matches the paper's own epistemic status: \texttt{thm:dual}(ii)/\texttt{cor:tri}
rest on computer-algebra verification plus the REU induction) and prove everything
below sorry-free from it.

\section*{Step 2: block evaluation (exact, finite algebra)}

Let $\eta^0$ be the block: bound pairs (state $3$) on $\{-L,\dots,0\}$, empty on
$\{1,\dots,L\}$. For a two-particle dual $\xi$ at positions $(x_1,x_2)$:
$N^-_{x-1}(\xi_i)=0$ (one particle per species), and on the block
$N^+_{x+1}(\eta^0_i)=\#\{y>x:\ y\le0,\ y\ge-L\}=(-x)\vee0$ for $x\ge-L$. Hence for
$x_1,x_2\in\{-L,\dots,0\}$ each exponent is $x_i+(-x_i)=0$ and
\[
  D^{\mathrm{tri}}(\xi,\eta^0)=\bone\{x_1\le0\}\,\bone\{x_2\le0\}
  \qquad(\text{i.e.\ }k=0\text{ in the paper's }q^{2k}\text{ with this block
  convention; prove the identity in this normalised form}).
\]
(If the encoding places the block boundary at a general site $b_0$, the constant is
$q^{2k}$ with $k=2b_0$; state whichever matches the chosen convention, with the
constant explicit.)

\section*{Step 3: the crossbridge identity}

Combining Steps 1--2 with the evaluation of $D^{\mathrm{tri}}$ at the initial dual
$\xi^{\mathrm{bp},a}$ (one particle of each species at site $a$):
$D^{\mathrm{tri}}(\xi^{\mathrm{bp},a},\eta)
=\eta_{1,a}\eta_{2,a}\,q^{2\left(a+N^+_{a+1}(\eta_1)\right)}
q^{2\left(a+N^+_{a+1}(\eta_2)\right)}\cdot q^{-2a-2a}\cdots$ --- state cleanly as:
\begin{equation}\label{eq:cb}
  \E_{\eta^0}\Big[\eta_{1,a}(s)\,\eta_{2,a}(s)\,
  q^{2\left(a+N^+_{a+1}(\eta_1(s))\right)+2\left(a+N^+_{a+1}(\eta_2(s))\right)}\Big]
  \;=\;q^{2k}\,\mathbb P_{(a,a)}\big(X_1(s)\le0,\ X_2(s)\le0\big),
\end{equation}
the left side defined directly from the process semigroup (no opaque objects
needed at finite $L$), the right side from the two-particle dual semigroup, and
$k$ the explicit block constant of Step 2. This is
\texttt{eq:crossbridge} of the paper at finite $L$ --- the identity the paper
verified to machine precision for $L\le6$; here it becomes a theorem for all $L$.

\begin{remark}[interface]
Do not modify the statement of \texttt{lem\_crossbridge} in the Crossover file; its
\texttt{sorry} continues to stand for the continuum embedding. Add a docstring
noting the finite-lattice core is proved in the new file. If the two-particle dual
in \texttt{DualPairWitness} uses a different but equivalent encoding of the pair
process, prove the bridge lemma or keep the new file's dual self-contained ---
whichever is cleaner.
\end{remark}

\end{document}
